\section{Literature review}

\begin{tcolorbox}[colback=black!5!white,colframe=black!75!black,
  title=About this chapter]
  This chapter provides a comprehensive review of the existing literature and foundational concepts relevant to the development of Type3 Asset Administration Shells (AAS) and their integration with Multi-Agent Systems (MAS) using the Belief-Desire-Intention (BDI) approach. It covers the basic principles of AAS, different types of AAS with a focus on Type3, and explores the architecture and functionalities of MAS, particularly in industrial applications. The chapter also delves into the BDI model, highlighting its significance in enabling autonomous decision-making within agent-based systems.
  
\end{tcolorbox}

\subsection{Asset Administration Shells(AAS)}
\begin{tcolorbox}[colback=black!5!white,colframe=black!75!black,
  title=About AAS]
  In this section , we explore the fundamental concepts of Asset Administration Shells (AAS), which serve as digital representations of physical assets in the context of Industry 4.0. We discuss the structure and components of AAS, including submodels and data types, and provide an overview of the different types of AAS, with a particular emphasis on Type3 AAS. This section aims to establish a solid understanding of AAS as a foundation for their integration with Multi-Agent Systems (MAS) in subsequent sections.
\end{tcolorbox}

\subsubsection{Basic concept of AAS}
\begin{tcolorbox}[colback=black!5!white,colframe=black!75!black,
  title=About AAS Concept]
  In this subsection, we delve into the basic concept of Asset Administration Shells (AAS), which are pivotal in the digital transformation of industrial assets. We discuss the role of AAS in providing a standardized digital representation of physical assets, enabling interoperability and communication within Industry 4.0 ecosystems. The subsection outlines the key features and functionalities of AAS, highlighting their significance in facilitating asset management, data exchange, and integration with other digital systems.
\end{tcolorbox}

\subsubsection{Submodels and data types}
\begin{tcolorbox}[colback=black!5!white,colframe=black!75!black,
  title=About Submodels and Data Types]
  This subsection focuses on the structure of Asset Administration Shells (AAS), specifically the concept of submodels and data types. We explore how submodels serve as modular components within the AAS, encapsulating specific functionalities and data related to the asset. The subsection also discusses various data types used in AAS, emphasizing their role in ensuring accurate representation and interoperability of asset information. This foundational understanding is crucial for comprehending the complexities of AAS and their integration with Multi-Agent Systems (MAS).
\end{tcolorbox}

\subsubsection{Types of AAS}
\begin{tcolorbox}[colback=black!5!white,colframe=black!75!black,
  title=About Types of AAS]
  In this subsection, we provide an overview of the different types of Asset Administration Shells (AAS) as defined in the literature. We categorize AAS into various types based on their functionalities, capabilities, and levels of integration with physical assets. The subsection highlights the distinctions between these types, with a particular focus on Type3 AAS, which is characterized by its advanced features and proactive capabilities. This classification serves as a basis for understanding the specific requirements and design considerations for implementing Type3 AAS in industrial applications.
\end{tcolorbox}
\subsubsection{AAS Type3}
\begin{tcolorbox}[colback=black!5!white,colframe=black!75!black,
  title=About AAS Type3]
  This subsection delves into the specifics of Type3 Asset Administration Shells (AAS), which represent an advanced category of AAS with enhanced functionalities. We discuss the defining characteristics of Type3 AAS, including their ability to support proactive behaviors, complex data management, and integration with intelligent systems. The subsection also explores the design principles and architectural considerations unique to Type3 AAS, highlighting their role in facilitating sophisticated asset management and interoperability within Industry 4.0 environments.
\end{tcolorbox}
\subsection{Multi Agent System (MAS)}
\begin{tcolorbox}[colback=black!5!white,colframe=black!75!black,
  title=About MAS]
  In this section, we explore the concept of Multi-Agent Systems (MAS) and their relevance in industrial applications. We discuss the fundamental principles of MAS, including agent autonomy, communication, and cooperation. The section also highlights various frameworks and architectures used in implementing MAS, emphasizing their significance in enhancing the capabilities of Asset Administration Shells (AAS) through intelligent agent-based interactions. This foundational understanding of MAS sets the stage for examining its integration with AAS in subsequent sections.
\end{tcolorbox}
\subsection{BDI (Belief-Desire-Intention) Approach}
\begin{tcolorbox}[colback=black!5!white,colframe=black!75!black,
  title=About BDI Approach]
  This section delves into the Belief-Desire-Intention (BDI) approach, a prominent model for designing intelligent agents within Multi-Agent Systems (MAS). We explore the core components of the BDI model, including beliefs, desires, and intentions, and discuss how these elements contribute to autonomous decision-making and goal-oriented behavior in agents. The section also highlights the advantages of using the BDI approach in industrial applications, particularly in enhancing the functionality and adaptability of Asset Administration Shells (AAS) through intelligent agent integration.
\end{tcolorbox}